% GENERATED FILE. Edit canonical lessons, then run npm run generate:books.
% canonical-source-hash: fnv1a64:e601d3f2c710b58f
% canonical-lessons: ML-C312-rogam, ML-C312-rogi, ML-C312-nalsu, ML-C312-kuttivayppu, ML-C312-ambulansu

\chapter{Disease, Patient, Nurse, Injection, Ambulance}
\label{ch:ml-disease-patient-nurse-injection-ambulance}
\hlchaptermodality{312}

\begin{chapteropening}
\textbf{By the end of this chapter:} \emph{I can say a disease, a patient, a nurse, an injection and an ambulance.}

Complete the last lesson of chapter 312: I can say a disease, a patient, a nurse, an injection and an ambulance.
\end{chapteropening}

\section[\texorpdfstring{\ml{രോഗം} (rōgaṁ)}{രോഗം (rōgaṁ)}]{\ml{രോഗം} (rōgaṁ) — a disease, an illness}

\label{lesson:ML-C312-rogam}



\begin{quote}
Before the new one: say the Malayalam for to bear fruit, then the Malayalam for to spread.
\end{quote}

\subsection*{You'll want to know: \ml{രോഗം}}
\textbf{\ml{രോഗം}} — \emph{rōgaṁ} — \textquotedblleft{}a disease, an illness\textquotedblright{}.

\begin{cousinweb}
\textbf{\ml{രോഗം}} (\emph{rōgaṁ}) $\leftarrow$ Sanskrit \emph{rōga}, \textquotedblleft{}disease\textquotedblright{}, taken over whole as a tatsama, with Malayalam's noun ending \emph{-ṁ} (CDIAL 10829).
\end{cousinweb}

\begin{grammarlens}[title={What you've built}]
One more word for this chapter.
\end{grammarlens}

\subsection*{Your turn}
\begin{itemize}
\raggedright
  \item \emph{Say it:} \emph{rōgaṁ}
  \item \emph{Say it:} \emph{rōgaṁ}, once more
  \item \emph{Say it:} say \emph{rōgaṁ}
  \item \emph{Recall:} say the Malayalam for to bear fruit, then the Malayalam for to spread, then say \emph{rōgaṁ} again
\end{itemize}

\begin{tcolorbox}[breakable,colback=teal!4,colframe=teal!35!black,title={Before you move on}]
What does \ml{രോഗം} mean? (\textquotedblleft{}A disease, an illness\textquotedblright{}.) Say it once more.
\end{tcolorbox}
\section[\texorpdfstring{\ml{രോഗി} (rōgi)}{രോഗി (rōgi)}]{\ml{രോഗി} (rōgi) — a patient}

\label{lesson:ML-C312-rogi}



\begin{quote}
Before the new one: say the Malayalam for to spread, then the Malayalam for a disease.
\end{quote}

\subsection*{You'll want to know: \ml{രോഗി}}
\textbf{\ml{രോഗി}} — \emph{rōgi} — \textquotedblleft{}a patient\textquotedblright{}.

\begin{cousinweb}
\textbf{\ml{രോഗി}} (\emph{rōgi}) $\leftarrow$ Sanskrit \emph{rōgin}, \textquotedblleft{}diseased, a sick person\textquotedblright{}, taken over whole as a tatsama.
\end{cousinweb}

\begin{grammarlens}[title={What you've built}]
One more word for this chapter.
\end{grammarlens}

\subsection*{Your turn}
\begin{itemize}
\raggedright
  \item \emph{Say it:} \emph{rōgi}
  \item \emph{Say it:} \emph{rōgi}, once more
  \item \emph{Say it:} say \emph{rōgi}
  \item \emph{Recall:} say the Malayalam for to spread, then the Malayalam for a disease, then say \emph{rōgi} again
\end{itemize}

\begin{tcolorbox}[breakable,colback=teal!4,colframe=teal!35!black,title={Before you move on}]
What does \ml{രോഗി} mean? (\textquotedblleft{}A patient\textquotedblright{}.) Say it once more.
\end{tcolorbox}
\section[\texorpdfstring{\ml{നഴ്സ്} (naḻsŭ)}{നഴ്സ് (naḻsŭ)}]{\ml{നഴ്സ്} (naḻsŭ) — a nurse}

\label{lesson:ML-C312-nalsu}



\begin{quote}
Before the new one: say the Malayalam for a disease, then the Malayalam for a patient.
\end{quote}

\subsection*{You'll want to know: \ml{നഴ്സ്}}
\textbf{\ml{നഴ്സ്}} — \emph{naḻsŭ} — \textquotedblleft{}a nurse\textquotedblright{}.

\begin{cousinweb}
\textbf{\ml{നഴ്സ്}} (\emph{naḻsŭ}) $\leftarrow$ English \emph{nurse}.
\end{cousinweb}

\begin{grammarlens}[title={What you've built}]
One more word for this chapter.
\end{grammarlens}

\subsection*{Your turn}
\begin{itemize}
\raggedright
  \item \emph{Say it:} \emph{naḻsŭ}
  \item \emph{Say it:} \emph{naḻsŭ}, once more
  \item \emph{Say it:} say \emph{naḻsŭ}
  \item \emph{Recall:} say the Malayalam for a disease, then the Malayalam for a patient, then say \emph{naḻsŭ} again
\end{itemize}

\begin{tcolorbox}[breakable,colback=teal!4,colframe=teal!35!black,title={Before you move on}]
What does \ml{നഴ്സ്} mean? (\textquotedblleft{}A nurse\textquotedblright{}.) Say it once more.
\end{tcolorbox}
\section[\texorpdfstring{\ml{കുത്തിവയ്പ്പ്} (kuttivayppŭ)}{കുത്തിവയ്പ്പ് (kuttivayppŭ)}]{\ml{കുത്തിവയ്പ്പ്} (kuttivayppŭ) — an injection}

\label{lesson:ML-C312-kuttivayppu}



\begin{quote}
Before the new one: say the Malayalam for a patient, then the Malayalam for a nurse.
\end{quote}

\subsection*{You'll want to know: \ml{കുത്തിവയ്പ്പ്}}
\textbf{\ml{കുത്തിവയ്പ്പ്}} — \emph{kuttivayppŭ} — \textquotedblleft{}an injection\textquotedblright{}.

\begin{cousinweb}
\textbf{\ml{കുത്തിവയ്പ്പ്}} (\emph{kuttivayppŭ}) is native Dravidian, \textquotedblleft{}a piercing-in\textquotedblright{}: \emph{kutti}, from the verb \textquotedblleft{}to pierce, prick\textquotedblright{} (DEDR 1719, with Tamil \emph{kuttu}), and \emph{vayppu}, \textquotedblleft{}a putting\textquotedblright{}, from the verb \textquotedblleft{}to place\textquotedblright{} (DEDR 5549).
\end{cousinweb}

\begin{grammarlens}[title={What you've built}]
One more word for this chapter.
\end{grammarlens}

\subsection*{Your turn}
\begin{itemize}
\raggedright
  \item \emph{Say it:} \emph{kuttivayppŭ}
  \item \emph{Say it:} \emph{kuttivayppŭ}, once more
  \item \emph{Say it:} say \emph{kuttivayppŭ}
  \item \emph{Recall:} say the Malayalam for a patient, then the Malayalam for a nurse, then say \emph{kuttivayppŭ} again
\end{itemize}

\begin{tcolorbox}[breakable,colback=teal!4,colframe=teal!35!black,title={Before you move on}]
What does \ml{കുത്തിവയ്പ്പ്} mean? (\textquotedblleft{}An injection\textquotedblright{}.) Say it once more.
\end{tcolorbox}
\section[\texorpdfstring{\ml{ആംബുലൻസ്} (āmbulansŭ)}{ആംബുലൻസ് (āmbulansŭ)}]{\ml{ആംബുലൻസ്} (āmbulansŭ) — an ambulance}

\label{lesson:ML-C312-ambulansu}



\begin{quote}
Before the new one: say the Malayalam for a nurse, then the Malayalam for an injection.
\end{quote}

\subsection*{You'll want to know: \ml{ആംബുലൻസ്}}
\textbf{\ml{ആംബുലൻസ്}} — \emph{āmbulansŭ} — \textquotedblleft{}an ambulance\textquotedblright{}.

\begin{cousinweb}
\textbf{\ml{ആംബുലൻസ്}} (\emph{āmbulansŭ}) $\leftarrow$ English \emph{ambulance}, from French \emph{hôpital ambulant}, \textquotedblleft{}walking hospital\textquotedblright{}.
\end{cousinweb}

\begin{grammarlens}[title={What you've built}]
One more word for this chapter.
\end{grammarlens}

\subsection*{Your turn}
\begin{itemize}
\raggedright
  \item \emph{Say it:} \emph{āmbulansŭ}
  \item \emph{Say it:} \emph{āmbulansŭ}, once more
  \item \emph{Say it:} say \emph{āmbulansŭ}
  \item \emph{Recall:} say the Malayalam for a nurse, then the Malayalam for an injection, then say \emph{āmbulansŭ} again
\end{itemize}

\begin{tcolorbox}[breakable,colback=teal!4,colframe=teal!35!black,title={Before you move on}]
What does \ml{ആംബുലൻസ്} mean? (\textquotedblleft{}An ambulance\textquotedblright{}.) Say it once more.
\end{tcolorbox}
